\section{Related Work}\label{sec:related}
% Compressed to ~half a column 2026-07-24 per the Meeting-44 mandate [1:20:02]:
% approaches + differences only. The full five-theme survey (drafted 2026-07-12,
% incl. Dr. Zhang's per-theme FATE-contrast closers) moved VERBATIM to
% Appendix~\ref{app:related} — restore-at-camera-ready candidate.

Fairness interventions are conventionally grouped by where they act: on the
training data before training (pre-processing), inside the training
objective (in-processing), or on a trained model's outputs
(post-processing) \cite{barocas2023}. FATE is pre-processing for
demonstrations: it adapts two ideas from classification fairness to the
trajectory setting. From instance reweighing \cite{kamirancalders2012} it
takes the practice of changing how much weight each demonstration carries
during training; from conditional statistical parity \cite{corbettdavies2017} it
takes the principle that a disparity counts only after permitted factors,
here demand, are adjusted for. The closest transportation
neighbor adds a fairness penalty inside one ride-hailing demand predictor's
training loss \cite{zheng2023}; that correction lives and dies with the one
model, while the demonstrations stay biased for every other consumer, which
is the gap FATE closes one stage earlier by addressing the problem at the data 
level. Fair data generation synthesizes
fair datasets outright \cite{xu2018fairgan,vanbreugel2021decaf}; FATE
generates nothing and edits real trajectories under a fidelity-protecting bound. Imitation
learning for mobility recovers driver policies faithfully (conditional GAIL
and its explainable variants \cite{zhang2019cgail,zhang2022cgail,pan2020xgail}),
but this lineage optimizes how faithfully behavior is reproduced and does not
audit the fairness of the reproduced allocation; FATE edits the
demonstrations such models consume, and reuses the driver-identity
verification line \cite{ren2020stsiamese} as its frozen fidelity gate. The
editing machinery descends from bounded adversarial perturbation
\cite{goodfellow2015fgsm,kurakin2017ifgsm,hu2023stifgsm}, read
constructively as in algorithmic recourse
\cite{ustun2019recourse,wachter2018counterfactual}: a minimal bounded change
prescribed for benefit rather than attack. Finally, the leveling-down
critique of fairness constraints \cite{parfit1997,mittelstadt2024,zietlow2022}
sets the standard we hold a fairness gain to: service outcomes of the under-served must actually
improve (\S\ref{sec:leveling}). Appendix~\ref{app:related} surveys these five
lines in full.
